\chapter{Large Numbers}\label{ch:g4:numbers}

Thousands were the ceiling of \cref{ch:g3:numbers}; now the ceiling
rises to \emph{millions}. The trick that makes huge numbers readable
is to cut their digits into groups of three --- the \omterm{def:g4:numbers:classes}{classes}.

\section{Classes of three digits}

\begin{definition}[Classes]\label{def:g4:numbers:classes}
Starting from the right, the digits of a number are grouped in
\emph{classes} of three: the class of \emph{units}, the class of
\emph{thousands}, the class of \emph{millions}. Each class has its
own units, tens and hundreds. For instance
\[
6\,309\,452 = 6 \text{ millions} + 309 \text{ thousands}
+ 452 \text{ units},
\]
read ``six million three hundred nine thousand four hundred
fifty-two''.
\end{definition}

\begin{omfigure}
\begin{tikzpicture}[scale=0.86]
  \draw (0,0) grid[xstep=1.15, ystep=0.75] (10.35,1.5);
  \node[font=\footnotesize, omThm] at (1.725,1.85) {millions};
  \node[font=\footnotesize, omDef] at (5.175,1.85) {thousands};
  \node[font=\footnotesize, omMeth] at (8.625,1.85) {units};
  \foreach \x/\l in {0.575/h, 1.725/t, 2.875/u,
                     4.025/h, 5.175/t, 6.325/u,
                     7.475/h, 8.625/t, 9.775/u}
    \node[font=\footnotesize] at (\x,1.12) {\l};
  \foreach \x/\d in {2.875/6, 4.025/3, 5.175/0, 6.325/9,
                     7.475/4, 8.625/5, 9.775/2}
    \node[font=\large] at (\x,0.37) {$\d$};
  \draw[omThm, thick] (0,0) rectangle (3.45,1.5);
  \draw[omDef, thick] (3.45,0) rectangle (6.9,1.5);
  \draw[omMeth, thick] (6.9,0) rectangle (10.35,1.5);
\end{tikzpicture}

{\small The class table of $6\,309\,452$ (h $=$ hundreds, t $=$
tens, u $=$ units of each class). Reading a large number means
reading each class, followed by its name.}
\end{omfigure}

\begin{example}[Writing with the classes]\label{ex:g4:numbers:write}
``Twelve million forty thousand five hundred'' has $12$ in the
millions class, $040$ in the thousands class, $500$ in the units
class:
\[
12\,040\,500 .
\]
The \omterm{def:g1:counting:zero}{zeros} of $040$ are essential: each class (except the leftmost)
must show its three digits.
\end{example}

\begin{example}[How many thousands?]\label{ex:g4:numbers:contained}
The number $6\,309\,452$ \emph{contains} $6\,309$ thousands
(everything left of the units class): $6\,309\,452 = 6\,309 \times
1\,000 + 452$. Asking ``how many thousands'' is not asking for one
digit, but for a whole prefix of the number.
\end{example}

\section{Comparing and rounding}

\begin{method}[Comparing large numbers]\label{met:g4:numbers:compare}
Exactly as in \cref{met:g3:numbers:compare}, with the \omterm{def:g4:numbers:classes}{classes} as a
shortcut: more digits wins; same length, compare from the left ---
class by class, then digit by digit.
\end{method}

\begin{example}\label{ex:g4:numbers:compare}
$2\,145\,890$ and $2\,148\,540$: millions equal ($2$), thousands
class $145 < 148$, so
\[
2\,145\,890 < 2\,148\,540 .
\]
\end{example}

\begin{definition}[Rounding]\label{def:g4:numbers:round}
To \emph{round}\index{rounding} a number to the nearest thousand
(or ten, hundred, million\,\dots), replace it by the closest multiple
of $1\,000$. If the number sits exactly halfway, round up. In
practice: look at the digit just below the rounding place --- $5$ or
more rounds up, $4$ or less rounds down.
\end{definition}

\begin{example}\label{ex:g4:numbers:round}
$37\,614$ rounded to the nearest thousand: the hundreds digit is
$6$, so \omterm{def:g4:numbers:round}{round} up: $38\,000$. Rounded to the nearest hundred:
$37\,600$. Rounded to the nearest ten thousand: $40\,000$.
\end{example}

\begin{omfigure}
\begin{tikzpicture}[scale=1.05]
  \draw[-Stealth] (-0.3,0) -- (10.8,0);
  \draw (0,-0.1) -- (0,0.1);
  \node[below, font=\small] at (0,-0.12) {$37\,000$};
  \draw (10,-0.1) -- (10,0.1);
  \node[below, font=\small] at (10,-0.12) {$38\,000$};
  \draw (5,-0.08) -- (5,0.08);
  \node[below, font=\footnotesize] at (5,-0.12) {$37\,500$};
  \fill[omThm] (6.14,0) circle (2pt)
    node[above, font=\small] {$37\,614$};
  \draw[-Stealth, omDef, thick] (6.14,0.42) .. controls (8,0.85) ..
    (9.9,0.28);
\end{tikzpicture}

{\small \omterm{def:g4:numbers:round}{Rounding} $37\,614$ to the nearest thousand: it lies past the
halfway mark $37\,500$, so it \omterm{def:g4:numbers:round}{rounds} up to $38\,000$.}
\end{omfigure}

\begin{example}[Orders of magnitude in the world]\label{ex:g4:numbers:world}
\omterm{def:g4:numbers:round}{Rounding} is how large numbers are actually used: a country of
$67\,842\,591$ inhabitants is ``about $68$ million''; a city of
$214\,376$ people is ``about $200\,000$''. The exact figure changes
every day anyway --- the rounded one carries the useful
information.
\end{example}

\section{Exercises}

\begin{exercise}[$\star$]\label{exo:g4:numbers:1}
Write in digits, with the \omterm{def:g4:numbers:classes}{classes} separated: ``three million five
hundred twenty thousand''; ``seven million eight''; ``two hundred
forty thousand sixty''.
\end{exercise}

\begin{exercise}[$\star$]\label{exo:g4:numbers:2}
Write in words: $5\,204\,300$; \quad $10\,050\,000$; \quad
$999\,999$.
\end{exercise}

\begin{exercise}[$\star$]\label{exo:g4:numbers:3}
In $8\,472\,156$: which digit is the tens of thousands? The hundreds
of the millions class? How many \emph{thousands} does the number
contain (see \cref{ex:g4:numbers:contained})?
\end{exercise}

\begin{exercise}[$\star$]\label{exo:g4:numbers:4}
Copy and complete with $<$ or $>$:
\[
980\,000 \;?\; 1\,020\,000, \qquad
3\,456\,000 \;?\; 3\,465\,000, \qquad
7\,090\,050 \;?\; 7\,090\,500 .
\]
\end{exercise}

\begin{exercise}[$\star$]\label{exo:g4:numbers:5}
Order from smallest to biggest: $2\,300\,000$; \ $980\,500$; \
$2\,030\,000$; \ $2\,299\,999$.
\end{exercise}

\begin{exercise}[$\star$]\label{exo:g4:numbers:6}
Round $63\,782$: to the nearest ten; to the nearest hundred; to the
nearest thousand; to the nearest ten thousand.
\end{exercise}

\begin{exercise}[$\star$]\label{exo:g4:numbers:7}
Round to the nearest million: $4\,721\,000$; \quad $6\,499\,999$;
\quad $9\,500\,000$.
\end{exercise}

\begin{exercise}[$\star$]\label{exo:g4:numbers:8}
Count by steps: three steps of $1\,000$ from $998\,500$; two steps
of $100\,000$ from $950\,000$. What happens to the \omterm{def:g4:numbers:classes}{classes} as you
cross a million?
\end{exercise}

\begin{exercise}[$\star$]\label{exo:g4:numbers:9}
A stadium holds $81\,338$ spectators. A newspaper writes ``more than
$80\,000$ fans''; a radio says ``about $81\,000$''. Are both
right? To which place did each \omterm{def:g4:numbers:round}{round}?
\end{exercise}

\begin{exercise}[$\star\star$]\label{exo:g4:numbers:10}
Using each of the digits $1$ to $7$ exactly once, write the biggest
possible seven-digit number, the smallest one, and the one closest to
$4\,000\,000$.
\end{exercise}

\begin{exercise}[$\star\star$]\label{exo:g4:numbers:11}
What is the biggest whole number that \omterm{def:g4:numbers:round}{rounds} to $50\,000$ when
\omterm{def:g4:numbers:round}{rounding} to the nearest thousand? And the smallest?
\end{exercise}
